\section{Engineering Scope and Evidence}
\label{sec:scope-and-evidence}

The public project is defined by the engineering specification in
Table~\ref{tab:design-specification}. The original administrative mechanism used to
select the frequency and feed type is intentionally omitted; the resulting 915~MHz
inset-fed configuration is retained as the fixed design input.

\begin{table}[H]
\centering
\caption{Public design specification.}
\label{tab:design-specification}
\begin{tabularx}{\textwidth}{@{}p{0.17\textwidth}p{0.33\textwidth}X@{}}
\toprule
\textbf{Category} & \textbf{Parameter} & \textbf{Specified value} \\
\midrule
Electrical & Design frequency & 915~MHz \\
Electrical & Feed topology & Inset-fed microstrip line \\
Material & Substrate & Rogers RT/duroid 5880 \\
Material & Relative permittivity & 2.20 \\
Material & Loss tangent & 0.0009 \\
Geometry & Substrate thickness & 1.575~mm \\
Geometry & Metal thickness used geometrically & 0.035~mm \\
Simulation & Frequency range & 0.6--1.2~GHz \\
Simulation & Boundaries & Open (add space), all sides \\
Simulation & Monitors & E field, H field, and far field at 0.915~GHz \\
\bottomrule
\end{tabularx}
\end{table}

\subsection{Evidence Classes}

\textbf{Analytical evidence} consists of the equations and independently executable
script used to regenerate the initial patch dimensions and inset-depth estimate.

\textbf{Simulated evidence} consists of CST screenshots. The exact $S_{11}$ minimum is
supported by a marker in the preserved plot. Directivity, main-beam direction,
beamwidth, and displayed efficiency values are read from the CST far-field information
panels.

\textbf{Measured evidence} is absent. No fabricated antenna, VNA sweep, calibrated
gain measurement, radiation-pattern measurement, or material characterization is
included.

\subsection{Public Claim Boundary}

The report claims only that the documented analytical model and preserved CST evidence
support the stated simulated behavior. It does not claim manufacturing readiness,
broadband performance, tolerance robustness, regulatory compliance, or measured
hardware performance.
